% ECONOMICS_PROBLEM: {"id": "C62-48", "title": "Sequential and recursive equilibrium with a distributional state", "jel_code": "C62", "source_id": "OP-0583"}
% !TeX program = xelatex
\documentclass[11pt,a4paper]{article}
\usepackage[margin=23mm]{geometry}
\usepackage{fontspec}
\setmainfont{Latin Modern Roman}
\usepackage{amsmath,amssymb,mathtools}
\usepackage{xcolor,enumitem,booktabs,longtable,array,xurl,hyperref}
\definecolor{muted}{HTML}{53636C}
\hypersetup{colorlinks=true,urlcolor=blue,linkcolor=blue}
\setlength{\parindent}{0pt}
\setlength{\parskip}{5pt}
\newcommand{\R}{\mathbb R}
\newcommand{\E}{\mathbb E}
\newcommand{\N}{\mathbb N}
\newcommand{\ind}{\mathbf 1}
\newcommand{\Var}{\operatorname{Var}}
\newcommand{\Cov}{\operatorname{Cov}}
\newcommand{\supp}{\operatorname{supp}}
\DeclareMathOperator*{\argmin}{arg\,min}
\DeclareMathOperator*{\argmax}{arg\,max}
\newcommand{\entrymeta}[3]{{\sffamily\small\color{muted}\textbf{\texttt{#1}}\quad|\quad #2\par #3\par}}
\newcommand{\Problemref}[1]{\texttt{#1}}
\newcommand{\JELref}[2]{\texttt{#2} (JEL \texttt{#1})}
\begin{document}
\section*{Sequential and recursive equilibrium with a distributional state}
\textbf{Problem ID:} \texttt{C62-48}\quad \textbf{JEL:} \texttt{C62}\par
% BEGIN ECONOMICS STATEMENT
\label{problem:OP-0583}
\label{beyond:MA-02}
\entrymeta{MA-02}{Heterogeneous-agent macroeconomics}{Published historical question and precise equivalence programme}
\subsection*{Definitions and assumptions}
Use the incomplete-market economy in MA-01. A sequential equilibrium specifies adapted prices and household consumption/savings plans at every aggregate and individual shock history. Households maximize expected discounted lifetime utility subject to their intertemporal budgets and borrowing bound, and markets clear at every relevant aggregate history. A recursive equilibrium uses the common aggregate state $s_t=(z_t,\mu_t)$ and individual state $(a_t,e_t)$, with the Bellman optimality and distribution-consistency conditions of MA-01.

A sequential equilibrium is representable on this reduced state if its prices, household policies, and conditional continuation values factor through these states: there exist measurable functions $R,w,g,V$ such that the corresponding sequential objects equal these functions at every reachable state, up to probability-null sets. The aggregate transition must also factor through $(z_t,\mu_t,z_{t+1})$. A recursive representation must specify an invariant state domain and consistent continuation choices, rather than only values along one realized path.
\subsection*{Formal question}
Characterize conditions under which a sequential equilibrium has such a recursive representation, and under which a recursive equilibrium generates a sequential optimum. Establish an equivalence theorem for a clearly specified equilibrium correspondence, including admissible initial distributions and any needed transversality or uniform-integrability conditions. Determine when auxiliary individual expected continuation values are redundant and when a continuation selection cannot be made as a function of $(z,\mu)$ alone.
\subsection*{Scope and status}
Equilibrium existence on the reduced state and representation of every sequential equilibrium are distinct demands. Histories with the same $(z,\mu)$ may support different continuation selections, so Markov dependence cannot simply be assumed. The question retains the entire wealth--efficiency distribution and concerns removal only of the auxiliary continuation coordinate. The source raises the relation as open in 2009; the present formulation separates the logical directions without asserting that all later work leaves them unresolved.
\subsection*{References for the source of the problem}
Heathcote, Storesletten, and Violante, \emph{Quantitative Macroeconomics with Heterogeneous Households}, NBER Working Paper~14768 (2009), \url{https://www.nber.org/system/files/working_papers/w14768/w14768.pdf}; verified published version, \emph{Annual Review of Economics} 1 (2009), Section~5.2, footnote~20, p.~342, \url{https://doi.org/10.1146/annurev.economics.050708.142922}. The footnote discusses Jianjun Miao, \emph{Competitive equilibria of economies with a continuum of consumers and aggregate shocks}, \emph{Journal of Economic Theory} 128(1) (2006), 274--298, as the earlier enlarged-state result.

\subsection*{Common conventions: Source mathematical conventions}

\textbf{Well-defined objects.}
All probability spaces and measurable variables are specified or inherited
from a local statistical model. Expectations used as finite targets require
the stated integrability. A decision rule or estimator is measurable with
respect to its available information; randomized rules may use an auxiliary
independent random variable. Norms without a subscript are Euclidean in
finite-dimensional spaces. An optimizer is requested only when attainment
is assured; otherwise the question concerns the optimal value and
$\varepsilon$-optimal admissible rules for every $\varepsilon>0$.

\textbf{Identification.}
A structural model $m\in\mathcal M$ implies an observable law
$P_m^{\rm obs}$ and target $\theta(m)$. The target is point identified on
$\mathcal M$ if
\[
 P_{m_1}^{\rm obs}=P_{m_2}^{\rm obs}
 \quad\Longrightarrow\quad\theta(m_1)=\theta(m_2)
 \qquad(m_1,m_2\in\mathcal M).
\]
When this fails, the sharp identified set is
\[
 \Theta(P;\mathcal M)=\{\theta(m):m\in\mathcal M,
                                  P_m^{\rm obs}=P\}.
\]
Specifying an intervention or estimand does not establish identification.
Positivity, exclusion, causal sufficiency, measurement restrictions, and
transport assumptions are substantive local inputs, never automatic
consequences of notation.

\textbf{Minimax risk and nontrivial inference.}
For a statistical experiment with data law $P^{(n)}$, target $\theta(P)$,
and nonnegative loss $\ell$, define
\[
 \mathcal R_n(\mathcal P,\ell)
   =\inf_{\widehat\theta_n}\sup_{P\in\mathcal P}
      \E_{P^{(n)}}\ell(\widehat\theta_n,\theta(P)).
\]
The infimum is over measurable estimators. A sharp rate is a sequence
$r_n$ for which $0<c\leq C<\infty$, independent of $n$, satisfies
$c r_n\leq\mathcal R_n\leq C r_n$ for all sufficiently large $n$.
Squared-error parametric risk is of order $n^{-1}$; the corresponding
root-mean-squared error is of order $n^{-1/2}$. These different conventions
are distinguished locally. A uniform asymptotic confidence region $C_n$
must satisfy
\[
 \liminf_{n\to\infty}\inf_{P\in\mathcal P}
     \Pr_{P^{(n)}}\{\theta(P)\in C_n\}\geq 1-\alpha,
 \qquad 0<\alpha<1,
\]
together with an informative diameter or expected-width requirement.
Coverage alone permits the vacuous whole parameter space. Testing questions
similarly impose both size control and power against a declared separated
alternative class.

\textbf{Binary causal baseline.}
When an entry explicitly invokes this baseline, one observes independent
copies of $O=(X,W,Y)$, with $W\in\{0,1\}$ and potential outcomes $Y(0),Y(1)$.
Assume consistency $Y=Y(W)$, no interference, conditional exchangeability
$(Y(0),Y(1))\perp W\mid X$, and overlap
$\epsilon\leq e(X)\leq1-\epsilon$ almost surely for a fixed
$0<\epsilon<1/2$, where
\[
 e(x)=\Pr(W=1\mid X=x),\qquad
 m_w(x)=\E[Y\mid W=w,X=x],\qquad
 \tau(x)=m_1(x)-m_0(x).
\]
These assumptions identify conditional mean effects almost everywhere
under the covariate law. Pointwise or uniform effects require appropriate
regularity and a specified support. Continuous treatment, longitudinal
data, adaptive experiments, and network interference require their own
assumptions in place of this baseline. Bounded outcomes, densities, and
smoothness are additional assumptions only where stated.

\textbf{Function classes.}
On $[0,1]^d$, $\mathcal H_d^s(L)$ is the isotropic Hölder ball of radius
$L>0$: put $k=\lceil s\rceil-1$ and $\gamma=s-k\in(0,1]$, bound derivatives
through order $k$, and require the order-$k$ derivatives to be
$\gamma$-Hölder with constant at most $L$. Thus integer orders use a
Lipschitz final derivative convention. The class
$\operatorname{BV}_d(L)$ consists of functions $f\in L^1((0,1)^d)$ whose
distributional derivative is a finite vector measure and for which
$\|f\|_{L^1}+|Df|((0,1)^d)\leq L$.
An $L^\infty$ bound or a particular pointwise version is imposed separately
when needed. Sparse and neural-network classes require a declared
dictionary or architecture and coefficient bounds.

An anisotropic H\"older class with declared block smoothness indices means
coordinatewise H\"older regularity in each block, uniformly in the other
blocks, with all corresponding derivative and H\"older bounds at most
the stated radius. Derivative orders and the integer-order convention in
each block are those just given. Mixed-derivative bounds are additional
restrictions only when explicitly stated.

\textbf{Decision and computational questions.}
Policy regret compares admissible feasible policies under the same law and
constraints. In computational entries, finite data and thresholds have
rational encodings; running time is measured in their total bit length.
Real-valued equilibrium search uses the explicitly declared algebraic or
FIXP encoding. An approximation ratio for maximization is written
$\operatorname{OPT}/\operatorname{ALG}$ and for minimization as
$\operatorname{ALG}/\operatorname{OPT}$, with zero optima and infeasible
instances handled locally. Historical approximation bounds are not
automatically current bounds.
% END ECONOMICS STATEMENT
\end{document}
